\section{Limitations and Threats to Validity}\label{sec:threats}

\textbf{Internal.} LLM output is non-deterministic even at fixed seeds on a hosted model, and results depend on the prompt and on the retry and token caps, which we tuned on the dev split only;
three seeds bound but do not remove this variance. The free tier limits us to 5 requests per minute and a daily token cap; if a run is resumed on a later day, the hosted model may have changed in between. The reduced seed grid (one seed for most baselines and ablations) leaves those rows without a variance estimate.
PIT results depend on the mutator set and on timeouts: PIT's default behaviour counts a mutant that makes the test time out as killed.

\textbf{External.} Three Java projects (Commons-cli on JUnit~5, Gson and Jsoup on JUnit~4) and 150 held-out tests, with smells restricted to the tsDetect set, do not support claims about other languages, test frameworks or larger suites.
Only one model family (\texttt{gpt-oss-120b}) was used, so the comparison between systems does not show how the gaps behave with a weaker or stronger model.

\textbf{Construct.} The PIT kill-set is a proxy for test strength: the same mutants can be killed by a weaker assertion, and mutants outside the sampled mutators are invisible. tsDetect and JNose disagree sharply on some smells
(ExceptionHandling, IgnoredTest), so ``smell removed'' depends on the detector rule; we report rules~i and~ii, and rule~iii only if the manual sample was completed. Empty baseline kill-sets are reported as unknown rather than safe. JaCoCo coverage, listed as a Phase~1 tool, is not used, so test strength is measured by the kill-set only.
Our evaluation set is 484 tests, not the 593 of~\cite{utrefactor}, because only those match the pinned project commits verbatim.

\section{Conclusion}\label{sec:conclusion}

We built a verified, agentic pipeline that accepts a test-smell refactoring only when it compiles, passes, removes the targeted smells and keeps every mutant the original test killed.
\textbf{G1:} \textbf{[TODO one sentence with the E3 vs E5-G1 number from results/E5]}. \textbf{G2:} \textbf{[TODO detector-consensus finding from results/E4 and E5-G2]}.
\textbf{G6:} \textbf{[TODO s3 vs b1 and b2 from results/E2, E3]}. Architectural lessons: tools, not the model, should decide acceptance; the baseline kill matrix must be computed before the patch is written
(a Phase~1 bug measured the refactored code instead); and failures must be captured per test so one error does not stop a run.
Future work: more projects and models, JaCoCo-style coverage next to the kill-set, parallel workers per project, and a learned planner trained on the cross-test memory.
